\section*{ID9215: Definition of Hm(Ω,t)}
\paragraph{Metadata.}
PiID: 1865436480229 | RTEID: RTE:P35783.8914:T2.4122 | UUID: 474c4d4c-32ec-553d-86cf-a3efb58da83f

\paragraph{Canonical Statement.}
rovides a basis for defining Berry and Sagnac phase contributions. Falsifiability / Test Handle: Experimentally construct a rotating double helix waveguide and measure curvature and torsion; confirm that timedependent modulations follow the predicted rotation frequency. 2.0.155 ID 304: Helical Transfer Function with Berry and Sagnac Phases Type: Operator Formal Statement: For orbital angular momentum (OAM) mode mmm at rotation frequency Ω\textbackslash{}OmegaΩ in a helical waveguide, the transfer function is \$Hm(Ω,t)=exp[i(ωmt+ΦmBerry+ΦmSagnac(Ω))]H\_m(\textbackslash{}Omega,t)=\textbackslash{}exp\textbackslash{}bigl[i(\textbackslash{}omega\_m\$ t + \textbackslash{}Phiˆ\{\textbackslash{}text\{Berry\}\} m + \$\textbackslash{}Phiˆ\{\textbackslash{}text\{Sagnac\}\} m(\textbackslash{}Omega))\textbackslash{}bigr]Hm(Ω,t)=exp[i(ωmt+ΦmBerry+ΦmSagnac(Ω))],\$ where \$ωm=ckz2+(m/R)2+2mΩ/n\textbackslash{}omega\_m\$ = c\textbackslash{}sqrt\{k\_zˆ2 + (m/R)ˆ2 + \$2m\textbackslash{}Omega/n\}ωm\$ = \$ckz2\$ + \$(m/R)2\$ + \$2mΩ/n,ΦmBerry\$ = \$(2lnφ/π)mt\textbackslash{}PhiΘ\{\textbackslash{}text\{Berry\}\} m\$ = \$(2\textbackslash{}ln\textbackslash{}varphi/\textbackslash{}pi)mtΦmBerry\$ = \$(2lnφ/π)mt, andΦmSagnac(Ω)\$ = \$(4πΩR2/λ0c)m\textbackslash{}Phiˆ\{\textbackslash{}text\{Sagnac\}\} m(\textbackslash{}Omega)\$ = (4\textbackslash{}pi\textbackslash{}Omega Rˆ2/

\paragraph{Canonical Equation.}
\[
Hm(Ω,t)=exp[i(ωmt+ΦmBerry+ΦmSagnac(Ω))]H_m(\Omega,t)=\exp\bigl[i(\omega_m
\]

\paragraph{Dictionary.}
Definition of Hm(Ω,t) — Hm(Ω,t)=exp[i(ωmt+ΦmBerry+ΦmSagnac(Ω))]H\_m(Ω,t)=exp l[i(ω\_m

\paragraph{Encyclopedia.}
Definition of Hm(Ω,t) is an algebraic definition, written Hm(Ω,t)=exp[i(ωmt+ΦmBerry+ΦmSagnac(Ω))]H\_m(Ω,t)=exp l[i(ω\_m. It is an algebraic definition expressing the quantity directly in terms of the variables on its right-hand side. It is built from Ω, x, p, m, B, r, defining Hm(Ω,t). Within the Trawin Zero Point registry it serves as one element of the framework's mathematical narrative.
